\documentclass{article}
\usepackage[utf8]{inputenc}
\usepackage{amsmath,amssymb}
\usepackage[most]{tcolorbox}
\usepackage{geometry}
\geometry{margin=2.5cm}

\newcommand{\step}[1]{\par\noindent\hangindent=3.5em\hangafter=1%
  \makebox[3.5em][l]{\textbf{#1.}}}
\newcommand{\steptag}[1]{\hfill{\small\textsf{[#1]}}}

\newtcolorbox{proofbox}[1]{
  colback=gray!5, colframe=gray!60,
  fonttitle=\bfseries, title={#1},
  breakable, enhanced,
  left=4pt, right=4pt, top=4pt, bottom=4pt,
  before skip=10pt, after skip=10pt
}

\begin{document}

\begin{proofbox}{Genuine-mass disintegration: for an exact signed idempote...}
\step{1} Genuine-mass disintegration: for an exact signed idempotent P with 0 < delta(P) <= 1/4, nonempty visible set W(P), height H, halo width a > 0 with G\_a = \{j : dist\_1(p\_j, conv W) > a*tau\} nonempty (tau = sqrt(delta)), fix for every row a vertex representation p\_j = sum\_v lambda\_jv p\_v over geometrically distinct row vertices; then every row index i satisfies g\_i <= M\_i\textasciicircum{}a + sum\_\{j in G\_a\} P\_ij\textasciicircum{}+ * (H - d\_j)/(H - a*tau), where g = P*1\_\{G\_a\}, d\_j = dist\_1(p\_j, conv W), P\_ij\textasciicircum{}+ = max(P\_ij, 0), and M\_i\textasciicircum{}a = sum\_\{j in G\_a\} P\_ij\textasciicircum{}+ * sum\_\{v : d\_v > a*tau\} lambda\_jv is positive mass supported entirely on HIDDEN row vertices at depth in (a*tau, H]. \steptag{validated} \\[6pt]
\step{1.1} Let C=conv\{p\_w : w in W(P)\} and d\_j=dist\_1(p\_j,C). Because G\_a is nonempty, some j has d\_j>a*tau, while the height definition gives d\_j<=H for every row; hence H>a*tau and H-a*tau>0. For any row vertex v with d\_v>a*tau, v is not visible since visible vertices lie in C and have distance 0; and d\_v<=H, so such vertices are hidden and have depth in (a*tau,H]. \steptag{validated} \\[4pt]
\step{1.2} For every row j, its fixed vertex representation p\_j=sum\_v lambda\_jv p\_v is a convex combination over geometrically distinct row vertices. Applying lem-residual-upper to C, the positive weights b\_v=lambda\_jv, no negative terms, m=sum\_v lambda\_jv=1, and q=p\_j gives the depth-convexity bound d\_j<=sum\_v lambda\_jv d\_v. \steptag{validated} \\[6pt]
\step{1.2.1} For a fixed row j, the phrase vertex representation in the root hypothesis means a finite convex representation over the chosen geometrically distinct row vertices: lambda\_jv>=0, sum\_v lambda\_jv=1, and p\_j=sum\_v lambda\_jv p\_v. The set C=conv W(P) is the convex hull of finitely many row points. \steptag{validated} \\[4pt]
\step{1.2.2} Use lem-residual-upper with C as above, positive coefficients b\_v=lambda\_jv, no negative coefficients c\_k, m=sum\_v lambda\_jv=1, q=p\_j, and p\_v the row vertices in the representation. Its conclusion is dist\_1(p\_j,C)<=sum\_v lambda\_jv dist\_1(p\_v,C), i.e. d\_j<=sum\_v lambda\_jv d\_v. \steptag{validated} \\[4pt]
\step{1.3} For every j in G\_a, define L\_j=sum\_\{v:d\_v<=a*tau\} lambda\_jv and U\_j=sum\_\{v:d\_v>a*tau\} lambda\_jv. Then L\_j+U\_j=1 and the previous bound plus d\_v<=a*tau on the L\_j part and d\_v<=H on the U\_j part gives d\_j<=a*tau*L\_j+H*U\_j=H-(H-a*tau)L\_j; since H-a*tau>0, L\_j<=(H-d\_j)/(H-a*tau). \steptag{validated} \\[6pt]
\step{1.3.1} For j in G\_a, the sets of vertices with d\_v<=a*tau and d\_v>a*tau partition the vertex representation, so L\_j+U\_j=1. Combining node 1.2 with d\_v<=a*tau on the first part and d\_v<=H on the second part gives d\_j<=a*tau*L\_j+H*U\_j. \steptag{validated} \\[6pt]
\step{1.3.1.1} Dependency bridge for the depth-convexity premise: for the fixed j in G\_a under discussion, validated node 1.2.2 applies to the fixed vertex representation p\_j=sum\_v lambda\_jv p\_v and gives d\_j <= sum\_v lambda\_jv d\_v. This supplies the bound that the parent previously cited as node 1.2, now through an explicit validated dependency. \steptag{validated} \\[4pt]
\step{1.3.1.2} Partition bridge: validated node 1.2.1 records that the fixed vertex representation has lambda\_jv >= 0 and sum\_v lambda\_jv = 1. The two conditions d\_v <= a*tau and d\_v > a*tau are complementary for each represented vertex v, so L\_j=sum\_\{v:d\_v<=a*tau\}lambda\_jv and U\_j=sum\_\{v:d\_v>a*tau\}lambda\_jv satisfy L\_j+U\_j=sum\_v lambda\_jv=1. \steptag{validated} \\[4pt]
\step{1.3.1.3} Combination bridge: split sum\_v lambda\_jv d\_v over the two complementary classes. On the class d\_v <= a*tau, nonnegativity of lambda\_jv gives sum\_\{d\_v<=a*tau\} lambda\_jv d\_v <= a*tau*L\_j. On the class d\_v > a*tau, validated node 1.1 gives d\_v <= H for every row vertex v, hence sum\_\{d\_v>a*tau\} lambda\_jv d\_v <= H*U\_j. Together with node 1.3.1.1, this gives d\_j <= sum\_v lambda\_jv d\_v <= a*tau*L\_j + H*U\_j, while node 1.3.1.2 gives L\_j+U\_j=1. \steptag{validated} \\[4pt]
\step{1.3.2} Substitute U\_j=1-L\_j to get d\_j<=H-(H-a*tau)L\_j. Node 1.1 gives H-a*tau>0, so rearranging yields L\_j<=(H-d\_j)/(H-a*tau). \steptag{validated} \\[6pt]
\step{1.3.2.1} Dependency bridge for the algebra step: for the fixed j in G\_a, node 1.3.1 supplies L\_j+U\_j=1 and d\_j<=a*tau*L\_j+H*U\_j. The equality L\_j+U\_j=1 gives U\_j=1-L\_j. \steptag{validated} \\[4pt]
\step{1.3.2.2} Algebraic rearrangement: using node 1.3.2.1, substitute U\_j=1-L\_j in d\_j<=a*tau*L\_j+H*U\_j to get d\_j<=a*tau*L\_j+H*(1-L\_j)=H-(H-a*tau)*L\_j. Validated node 1.1 gives H-a*tau>0, so this inequality implies (H-a*tau)*L\_j<=H-d\_j and therefore L\_j<=(H-d\_j)/(H-a*tau). \steptag{validated} \\[4pt]
\step{1.4} For every row index i, g\_i=(P*1\_\{G\_a\})\_i=sum\_\{j in G\_a\} P\_ij<=sum\_\{j in G\_a\} P\_ij\textasciicircum{}+. Splitting 1=L\_j+U\_j in each positive coefficient gives sum\_\{j in G\_a\} P\_ij\textasciicircum{}+=M\_i\textasciicircum{}a+sum\_\{j in G\_a\} P\_ij\textasciicircum{}+ L\_j, where M\_i\textasciicircum{}a=sum\_\{j in G\_a\} P\_ij\textasciicircum{}+ U\_j. Multiplying the bound on L\_j by P\_ij\textasciicircum{}+>=0 and summing over j in G\_a yields g\_i<=M\_i\textasciicircum{}a+sum\_\{j in G\_a\} P\_ij\textasciicircum{}+*(H-d\_j)/(H-a*tau), and the first step proves that M\_i\textasciicircum{}a is positive mass supported entirely on hidden row vertices at depth in (a*tau,H]. \steptag{validated} \\[6pt]
\step{1.4.1} For each i, the definition g=P*1\_\{G\_a\} gives g\_i=sum\_\{j in G\_a\} P\_ij. Since P\_ij<=P\_ij\textasciicircum{}+=max(P\_ij,0) for each j, g\_i<=sum\_\{j in G\_a\} P\_ij\textasciicircum{}+. \steptag{validated} \\[4pt]
\step{1.4.2} For fixed row index i and j in G\_a, set L\_j=sum\_\{v:d\_v<=a*tau\} lambda\_jv and U\_j=sum\_\{v:d\_v>a*tau\} lambda\_jv. The needed shallow-mass estimate L\_j <= (H-d\_j)/(H-a*tau) is proved locally from the fixed vertex representation, the residual-upper depth-convexity bound, and H-a*tau>0, rather than imported from pending node 1.3. Since P\_ij\textasciicircum{}+>=0, summing this estimate over j in G\_a gives sum\_\{j in G\_a\} P\_ij\textasciicircum{}+ L\_j <= sum\_\{j in G\_a\} P\_ij\textasciicircum{}+*(H-d\_j)/(H-a*tau). Also L\_j+U\_j=1, so sum\_\{j in G\_a\} P\_ij\textasciicircum{}+ = sum\_\{j in G\_a\} P\_ij\textasciicircum{}+ U\_j + sum\_\{j in G\_a\} P\_ij\textasciicircum{}+ L\_j, and the first term is exactly M\_i\textasciicircum{}a. Together with node 1.4.1 this gives the displayed bound; validated node 1.1 gives the claimed hidden-depth support of M\_i\textasciicircum{}a. \steptag{validated} \\[6pt]
\step{1.4.2.1} Local shallow-mass bridge: fix j in G\_a. Validated node 1.2.1 gives lambda\_jv>=0, sum\_v lambda\_jv=1, and the two complementary classes d\_v<=a*tau and d\_v>a*tau, so L\_j+U\_j=1. Validated node 1.2.2 gives d\_j<=sum\_v lambda\_jv d\_v. Splitting this sum, the d\_v<=a*tau part is at most a*tau*L\_j, and validated node 1.1 gives d\_v<=H for every represented row vertex, so the d\_v>a*tau part is at most H*U\_j. Hence d\_j<=a*tau*L\_j+H*U\_j. With U\_j=1-L\_j and H-a*tau>0 from validated node 1.1, this rearranges to L\_j<=(H-d\_j)/(H-a*tau). \steptag{validated} \\[4pt]
\step{1.4.2.2} Weighted summation bridge: for the fixed row i, P\_ij\textasciicircum{}+=max(P\_ij,0)>=0 for every j. Multiplying the estimate from node 1.4.2.1 by P\_ij\textasciicircum{}+ and summing over the finite set G\_a preserves the inequality, giving sum\_\{j in G\_a\} P\_ij\textasciicircum{}+ L\_j <= sum\_\{j in G\_a\} P\_ij\textasciicircum{}+*(H-d\_j)/(H-a*tau). \steptag{validated} \\[4pt]
\step{1.4.2.3} Assembly and support bridge: because L\_j+U\_j=1, sum\_\{j in G\_a\} P\_ij\textasciicircum{}+=sum\_\{j in G\_a\} P\_ij\textasciicircum{}+ U\_j+sum\_\{j in G\_a\} P\_ij\textasciicircum{}+ L\_j, and by the root definition of M\_i\textasciicircum{}a with U\_j=sum\_\{v:d\_v>a*tau\} lambda\_jv the first sum is exactly M\_i\textasciicircum{}a. Combining validated node 1.4.1 with node 1.4.2.2 yields g\_i<=M\_i\textasciicircum{}a+sum\_\{j in G\_a\} P\_ij\textasciicircum{}+*(H-d\_j)/(H-a*tau). The coefficients P\_ij\textasciicircum{}+ lambda\_jv in M\_i\textasciicircum{}a are nonnegative, and validated node 1.1 says every row vertex with d\_v>a*tau is hidden with depth in (a*tau,H], so M\_i\textasciicircum{}a is supported entirely on those hidden vertices. \steptag{validated} \\[4pt]
\end{proofbox}

\end{document}
